\chapter*{Preface}
\addcontentsline{toc}{chapter}{Preface}

Every book of this cycle has, at some point, gestured toward oscillation without making it the subject: Book 6 sketched paired curvature modes, Book 8 extended the flow to complex time, Book 9 raised the question of periodic geodesics only to defer it. The present book completes the geometric phase of the RSVP program by formalizing what all three were reaching for, the lamphron-lamphrodyne cycle, a coupled oscillation between local tension and global relaxation that sustains intelligibility across scales. Entropy never decreases, exactly as this series has maintained since Book 1, yet it oscillates locally, producing temporal coherence without expansion. This is RSVP's model of breathing comprehension: the universe as a self-oscillating manifold of meaning, and the last of this cycle's five books to draw on the others' deferred material rather than defer further material of its own.

\chapter*{Diagrammatic Overview}
\addcontentsline{toc}{chapter}{Diagrammatic Overview}

Lamphron and lamphrodyne name the two conjugate modes of a single manifold: lamphron concentrates curvature, a compressive phase in which local tension builds, and lamphrodyne releases it, an expansive phase in which that tension smooths back out. Every chapter of this book studies some aspect of the coupling between these two modes, from the differential equations governing their exchange to the ethical and cognitive consequences of a coupling that fails to complete.

\chapter*{Reading Note}
\addcontentsline{toc}{chapter}{Reading Note}

Treat every cycle in this book as a heartbeat of the plenum: not a decorative metaphor but the literal claim that compression without release, or release without compression, is not a milder version of the lamphron-lamphrodyne cycle but its failure, in exactly the sense a heartbeat that only contracts, or only relaxes, is not a gentler heartbeat but the absence of one.

\part{The Mathematics of Cyclic Curvature}

\chapter{Defining Lamphron and Lamphrodyne Modes}

The curvature tensor of Book 6 splits into dual components, $R_{ij} = R^{(+)}_{ij} + R^{(-)}_{ij}$, with $R^{(+)}$ the lamphron, compressive mode and $R^{(-)}$ the lamphrodyne, expansive mode. Coupling equations governing their exchange are
\[
\partial_t R^{(+)}_{ij} = \omega R^{(-)}_{ij}, \qquad \partial_t R^{(-)}_{ij} = -\omega R^{(+)}_{ij},
\]
with $\omega$ interpreted as a curvature-frequency parameter setting the rate of exchange between the two modes. This is the same linear structure that governs any pair of conjugate harmonic variables, and it carries a direct conservation law: total curvature energy, defined in Chapter 2 below, remains constant under this exchange, so that the cycle redistributes curvature between its two modes without any of it being created or destroyed by the coupling itself.

